% ============================================================================
% CAPÍTULO 3 - FINE-TUNING DO MODELO
% ============================================================================

\chapter{Fine-Tuning do Modelo}
\label{cap:fine-tuning}

Este capítulo descreve o processo de fine-tuning do modelo de linguagem para o domínio médico. O objetivo é adaptar um modelo generalista (Mistral 7B) para responder perguntas médicas em português com maior precisão e relevância clínica.

% ============================================================================
\section{Modelo Base: Mistral 7B Instruct}
\label{sec:modelo-base}

\subsection{Justificativa da Escolha}

Para o fine-tuning, escolhemos o \textbf{Mistral 7B Instruct v0.3} como modelo base. Esta decisão foi tomada após análise comparativa de alternativas, considerando os seguintes critérios:

\begin{table}[H]
\centering
\caption{Comparação de modelos candidatos para fine-tuning}
\label{tab:comparacao-modelos}
\begin{tabular}{lcccc}
\toprule
\textbf{Modelo} & \textbf{Parâmetros} & \textbf{Contexto} & \textbf{Licença} & \textbf{VRAM (4-bit)} \\
\midrule
Mistral 7B Instruct & 7.3B & 32K tokens & Apache 2.0 & ~5GB \\
LLaMA 2 7B & 7B & 4K tokens & Meta License & ~5GB \\
BioMistral 7B & 7.3B & 32K tokens & Apache 2.0 & ~5GB \\
Falcon 7B & 7B & 2K tokens & Apache 2.0 & ~5GB \\
\bottomrule
\end{tabular}
\end{table}

\begin{importantbox}[Por que não usar BioMistral?]
O BioMistral já possui conhecimento médico pré-treinado, o que dificultaria demonstrar o \textbf{delta} (melhoria) obtido com o fine-tuning. Ao usar o Mistral 7B generalista, conseguimos evidenciar claramente a diferença entre as respostas antes e depois do treinamento especializado.
\end{importantbox}

\subsection{Características do Mistral 7B}

O Mistral 7B apresenta características técnicas que o tornam ideal para este projeto:

\begin{itemize}
    \item \textbf{Arquitetura}: Transformer decoder-only com 32 camadas
    \item \textbf{Sliding Window Attention}: Janela de 4096 tokens com contexto efetivo de 32K
    \item \textbf{Grouped-Query Attention (GQA)}: Reduz uso de memória mantendo qualidade
    \item \textbf{Instruct-tuned}: Versão já otimizada para seguir instruções
    \item \textbf{Multilíngue}: Bom desempenho em português, mesmo sem treinamento específico
\end{itemize}

% ============================================================================
\section{Dataset: MedQuAD e Abordagem Híbrida}
\label{sec:dataset-medquad}

\subsection{Escolha do Dataset}

Para o fine-tuning, utilizamos o \textbf{MedQuAD} (Medical Question Answering Dataset), um dataset de perguntas e respostas médicas criado pelo National Institutes of Health (NIH) dos Estados Unidos, combinado com uma abordagem híbrida que inclui dados em inglês e conhecimento médico brasileiro.

\begin{table}[H]
\centering
\caption{Comparação de versões do MedQuAD disponíveis}
\label{tab:versoes-medquad-ft}
\begin{tabular}{lccc}
\toprule
\textbf{Dataset (HuggingFace)} & \textbf{Registros} & \textbf{Com Resposta} & \textbf{Escolhido} \\
\midrule
\texttt{lavita/MedQuAD} & 47.441 & $\sim$16.400 & \checkmark \\
\texttt{keivalya/MedQuad-...} & 16.400 & 16.400 & \\
\texttt{Tonic/medquad} & 15.549 & 15.549 & \\
\bottomrule
\end{tabular}
\end{table}

\begin{warningbox}[Limitação do Dataset]
O dataset \texttt{lavita/MedQuAD} possui 47.441 registros, porém \textbf{$\sim$65\% (31.034 registros) têm \texttt{answer=None}} devido a remoções por copyright solicitadas pelas instituições originais. O dataset utilizável contém aproximadamente \textbf{16.400 pares QA completos}.
\end{warningbox}

Optamos pela versão \texttt{lavita/MedQuAD} por oferecer:

\begin{itemize}
    \item \textbf{Metadados ricos}: códigos UMLS, tipos semânticos, sinônimos
    \item \textbf{Fontes documentadas}: NCI, GHR, MedlinePlus, NIDDK, NINDS, NIHSeniorHealth, NHLBI, CDC
    \item \textbf{Licença clara}: Creative Commons BY 4.0
\end{itemize}

\subsection{Fontes do MedQuAD}

O dataset original do MedQuAD compila informações de 8 fontes do NIH, cobrindo diversas especialidades médicas:

\begin{table}[H]
\centering
\caption{Fontes do dataset MedQuAD e cobertura temática}
\label{tab:fontes-medquad}
\begin{tabular}{llc}
\toprule
\textbf{Fonte} & \textbf{Cobertura} & \textbf{Pares QA (úteis)} \\
\midrule
National Cancer Institute (NCI) & Oncologia & $\sim$250 \\
Genetics Home Reference (GHR) & Genética & $\sim$5.400 \\
MedlinePlus Health Topics & Saúde geral & $\sim$980 \\
NIDDK & Diabetes, digestivo, renal & $\sim$1.100 \\
NINDS & Neurologia & $\sim$1.100 \\
NIHSeniorHealth & Geriatria & $\sim$750 \\
NHLBI & Cardiologia, pneumologia & $\sim$550 \\
CDC & Doenças infecciosas & $\sim$270 \\
\midrule
\textbf{Total (com resposta)} & & \textbf{$\sim$16.400} \\
\bottomrule
\end{tabular}
\end{table}

\subsection{Abordagem Híbrida: PT + EN + Brasileiro}
\label{subsec:abordagem-hibrida}

Com base em pesquisa sobre fine-tuning multilíngue, adotamos uma abordagem híbrida que combina três fontes de dados para maximizar a qualidade do modelo:

\begin{enumerate}
    \item \textbf{MedQuAD traduzido (78\%):} $\sim$16.400 pares traduzidos para português
    \item \textbf{MedQuAD original em inglês (20\%):} $\sim$4.100 pares preservados no idioma original
    \item \textbf{QA de medicina brasileira (2\%):} 55 pares sobre doenças endêmicas, SUS, medicamentos brasileiros
\end{enumerate}

\begin{infobox}[Fundamentação da Abordagem Híbrida]
Estudos recentes demonstram que: (1) dados traduzidos automaticamente introduzem ruído que o modelo aprende a imitar; (2) fine-tuning exclusivo em um idioma pode degradar capacidades multilíngues; (3) apenas 40 exemplos multilíngues preservam capacidades cross-lingual. Nossa proporção de 20\% em inglês mitiga esses problemas.
\end{infobox}

\subsection{Tradução para Português}
\label{subsec:traducao}

O dataset MedQuAD original está em inglês. Para treinar o modelo a responder em português brasileiro, implementamos um pipeline de tradução automatizada.

\subsubsection{Estratégia de Tradução}

Utilizamos a biblioteca \texttt{deep-translator} que oferece acesso gratuito à API do Google Translate:

\begin{lstlisting}[caption={Função de tradução com tratamento de textos longos}, label={lst:traducao}]
from deep_translator import GoogleTranslator

def translate_to_portuguese(text: str) -> str:
    """Traduz texto de ingles para portugues brasileiro."""
    if not text or len(text.strip()) == 0:
        return text

    translator = GoogleTranslator(source='en', target='pt')

    # Textos longos precisam ser divididos (limite de 5000 chars)
    if len(text) > 4500:
        chunks = [text[i:i+4500] for i in range(0, len(text), 4500)]
        translated_chunks = [translator.translate(chunk) for chunk in chunks]
        return ' '.join(translated_chunks)

    return translator.translate(text)
\end{lstlisting}

A tradução dos $\sim$16.400 registros leva aproximadamente \textbf{4--6 horas} devido ao volume de chamadas à API do Google Translate e ao \textit{rate limiting} imposto pelo serviço. O sistema de checkpoints descrito a seguir permite retomar o processo em caso de interrupção.

\begin{infobox}[Dataset já traduzido disponível no repositório]
Para evitar que cada usuário precise repetir o processo de tradução, o dataset final já traduzido e formatado está commitado no repositório em \texttt{data/processed/}:
\begin{itemize}
    \item \texttt{medquad\_pt.jsonl} --- MedQuAD traduzido para português ($\sim$16.400 pares)
    \item \texttt{medquad\_en\_sample.jsonl} --- Amostra em inglês ($\sim$4.100 pares)
    \item \texttt{brazilian\_medical\_qa.jsonl} --- QA de doenças brasileiras (55 pares)
    \item \texttt{training\_data.jsonl} --- Dataset final híbrido ($\sim$20.500 pares, pronto para treino)
\end{itemize}
O script \texttt{03\_preparar\_dataset.py} só precisa ser re-executado se o usuário desejar regenerar a tradução ou alterar a composição do dataset.
\end{infobox}

\subsubsection{Sistema de Checkpoints}

Dado o volume do dataset ($\sim$16.400 registros úteis), implementamos um sistema de checkpoints que salva o progresso a cada 1.000 traduções. Isso permite:

\begin{itemize}
    \item Retomar o processo em caso de interrupção
    \item Monitorar o progresso da tradução
    \item Evitar retrabalho em caso de erros de rede
\end{itemize}

\begin{lstlisting}[caption={Sistema de checkpoint para tradução}, label={lst:checkpoint}]
CHECKPOINT_INTERVAL = 1000

for i, row in enumerate(dataset):
    # Pular registros sem resposta (copyright)
    if row['answer'] is None:
        continue

    # Traduzir pergunta e resposta
    translated_question = translate_to_portuguese(row['question'])
    translated_answer = translate_to_portuguese(row['answer'])

    translated_data.append({
        "instruction": "Responda a seguinte pergunta medica...",
        "input": translated_question,
        "output": translated_answer
    })

    # Salvar checkpoint a cada 1000 registros
    if (i + 1) % CHECKPOINT_INTERVAL == 0:
        save_checkpoint(translated_data, i + 1)
\end{lstlisting}

\subsection{Formato Alpaca}
\label{subsec:formato-alpaca}

O dataset traduzido é formatado no padrão \textbf{Alpaca}, amplamente utilizado para instruction tuning de \acrshort{llm}s:

\begin{lstlisting}[caption={Estrutura do formato Alpaca}, label={lst:alpaca-format}]
{
    "instruction": "Responda a seguinte pergunta medica de forma
                    clara e detalhada.",
    "input": "Quais sao os sintomas do diabetes tipo 2?",
    "output": "Os sintomas mais comuns do diabetes tipo 2 incluem:
               aumento da sede (polidipsia), miccao frequente
               (poliuria), fadiga, visao turva, cicatrizacao lenta
               de feridas, formigamento nas maos e pes..."
}
\end{lstlisting}

O template de prompt utilizado durante o treinamento é:

\begin{tcolorbox}[colback=code-bg, colframe=fiap-gray, title=Template Alpaca em Português]
\ttfamily\small
Abaixo está uma instrução que descreve uma tarefa, junto com uma entrada que fornece contexto adicional. Escreva uma resposta que complete adequadamente a solicitação.\\[0.5em]
\#\#\# Instrução:\\
\{instruction\}\\[0.5em]
\#\#\# Entrada:\\
\{input\}\\[0.5em]
\#\#\# Resposta:\\
\{output\}
\end{tcolorbox}

O arquivo final é salvo no formato JSONL (JSON Lines), com um registro por linha:

\begin{lstlisting}[style=bash, caption={Exemplo do arquivo training\_data.jsonl}]
{"instruction": "Responda...", "input": "O que e hipertensao?", "output": "A hipertensao arterial..."}
{"instruction": "Responda...", "input": "Quais as causas do diabetes?", "output": "O diabetes tipo 1..."}
\end{lstlisting}

% ============================================================================
\section{Técnica: QLoRA}
\label{sec:qlora}

Para realizar o fine-tuning de forma eficiente em hardware limitado (Google Colab Free com GPU T4), utilizamos a técnica \textbf{QLoRA} (Quantized Low-Rank Adaptation).

\subsection{Fundamentos do LoRA}

O \acrfull{lora} é uma técnica de fine-tuning eficiente que, em vez de atualizar todos os parâmetros do modelo, treina apenas matrizes de baixo rank que são adicionadas às camadas existentes:

\begin{equation}
W' = W + \Delta W = W + BA
\end{equation}

Onde:
\begin{itemize}
    \item $W \in \mathbb{R}^{d \times k}$ é a matriz de pesos original (congelada)
    \item $B \in \mathbb{R}^{d \times r}$ e $A \in \mathbb{R}^{r \times k}$ são as matrizes treináveis
    \item $r \ll \min(d, k)$ é o \textit{rank} do LoRA
\end{itemize}

\subsection{Quantização 4-bit}

O QLoRA combina LoRA com quantização de 4 bits, reduzindo drasticamente o uso de memória:

\begin{table}[H]
\centering
\caption{Comparação de uso de memória por técnica de fine-tuning}
\label{tab:memoria-tecnicas}
\begin{tabular}{lcc}
\toprule
\textbf{Técnica} & \textbf{VRAM (Mistral 7B)} & \textbf{Parâmetros Treináveis} \\
\midrule
Full Fine-Tuning (FP16) & ~28 GB & 7.3B (100\%) \\
LoRA (FP16) & ~14 GB & ~42M (0.6\%) \\
\textbf{QLoRA (4-bit)} & \textbf{~5 GB} & \textbf{~42M (0.6\%)} \\
\bottomrule
\end{tabular}
\end{table}

A configuração de quantização utilizada:

\begin{lstlisting}[caption={Configuração de quantização 4-bit}, label={lst:bnb-config}]
from transformers import BitsAndBytesConfig

bnb_config = BitsAndBytesConfig(
    load_in_4bit=True,              # Quantizacao 4-bit
    bnb_4bit_quant_type="nf4",      # Normal Float 4-bit
    bnb_4bit_compute_dtype=torch.float16,  # Computacao em FP16
    bnb_4bit_use_double_quant=True  # Quantizacao dupla (economia extra)
)
\end{lstlisting}

\subsection{Configuração do LoRA}

Os hiperparâmetros do LoRA foram escolhidos para maximizar a capacidade de aprendizado dentro das limitações de hardware:

\begin{lstlisting}[caption={Configuração LoRA para o fine-tuning}, label={lst:lora-config}]
from peft import LoraConfig

lora_config = LoraConfig(
    r=64,                    # Rank (maior = mais capacidade)
    lora_alpha=128,          # Scaling factor (2x o rank)
    lora_dropout=0.05,       # Regularizacao
    bias="none",             # Nao treinar bias
    task_type="CAUSAL_LM",   # Modelagem de linguagem causal

    # Camadas a serem treinadas
    target_modules=[
        "q_proj", "k_proj", "v_proj", "o_proj",  # Attention
        "gate_proj", "up_proj", "down_proj"       # MLP
    ]
)
\end{lstlisting}

\begin{infobox}[Escolha do Rank]
O rank $r=64$ foi escolhido como um equilíbrio entre capacidade de aprendizado e uso de memória. Ranks maiores (128, 256) podem capturar mais nuances, mas aumentam significativamente o tempo de treino e o risco de overfitting. Para domínios especializados como medicina, $r \in [32, 64]$ é geralmente suficiente.
\end{infobox}

% ============================================================================
\section{Processo de Treino}
\label{sec:processo-treino}

\subsection{Ambiente de Execução}

O projeto disponibiliza \textbf{dois notebooks} para o fine-tuning, permitindo flexibilidade na escolha do ambiente:

\begin{table}[H]
\centering
\caption{Notebooks disponíveis para fine-tuning}
\label{tab:notebooks-finetuning}
\begin{tabular}{lp{5cm}l}
\toprule
\textbf{Notebook} & \textbf{Ambiente} & \textbf{Requisitos} \\
\midrule
\texttt{03\_fine\_tuning.ipynb} & Google Colab Free & Conta Google \\
\texttt{03\_fine\_tuning\_local.ipynb} & Servidor/PC local & GPU 16GB+ VRAM \\
\bottomrule
\end{tabular}
\end{table}

O notebook para \textbf{Google Colab} é otimizado para sessões limitadas (checkpoints frequentes, Google Drive). O notebook \textbf{local} auto-detecta a GPU disponível e ajusta batch size automaticamente, suportando também Apple Silicon (MPS).

Para o ambiente Colab Free, as especificações são:

\begin{table}[H]
\centering
\caption{Especificações do ambiente de treinamento}
\label{tab:ambiente-treino}
\begin{tabular}{ll}
\toprule
\textbf{Componente} & \textbf{Especificação} \\
\midrule
Plataforma & Google Colab Free \\
GPU & NVIDIA Tesla T4 \\
VRAM & 15 GB \\
RAM & 12.7 GB \\
Limite de sessão & ~12 horas \\
Armazenamento & Google Drive (15 GB free) \\
\bottomrule
\end{tabular}
\end{table}

\subsection{Hiperparâmetros de Treinamento}

Os hiperparâmetros foram configurados para otimizar a qualidade do modelo dentro das restrições de hardware e tempo:

\begin{lstlisting}[caption={Configuração completa do treinamento}, label={lst:training-args}]
from transformers import TrainingArguments

training_args = TrainingArguments(
    output_dir="./checkpoints",

    # Epochs e batching
    num_train_epochs=5,                    # 5 epochs para qualidade
    per_device_train_batch_size=2,         # Batch size (limite T4)
    gradient_accumulation_steps=8,         # Effective batch = 16

    # Learning rate
    learning_rate=2e-4,                    # Taxa padrao para LoRA
    warmup_ratio=0.03,                     # 3% warmup
    weight_decay=0.01,                     # Regularizacao L2
    lr_scheduler_type="cosine",            # Decay suave

    # Checkpoints (CRITICO para Colab Free!)
    save_strategy="steps",
    save_steps=500,                        # Salvar a cada ~30-40 min
    save_total_limit=3,                    # Manter 3 checkpoints

    # Avaliacao
    evaluation_strategy="steps",
    eval_steps=500,
    load_best_model_at_end=True,

    # Precisao
    fp16=True,                             # Mixed precision
)
\end{lstlisting}

\begin{table}[H]
\centering
\caption{Resumo dos hiperparâmetros de treinamento}
\label{tab:hiperparametros}
\begin{tabular}{lll}
\toprule
\textbf{Parâmetro} & \textbf{Valor} & \textbf{Justificativa} \\
\midrule
Epochs & 5 & Qualidade máxima, convergência completa \\
Batch size efetivo & 16 & $2 \times 8$ (grad accumulation) \\
Learning rate & $2 \times 10^{-4}$ & Padrão para LoRA \\
Warmup & 3\% & Estabilidade inicial \\
LR scheduler & Cosine & Decay suave até o final \\
Checkpoints & 500 steps & Recuperação de sessão \\
\bottomrule
\end{tabular}
\end{table}

\subsection{Sistema de Checkpoints para Colab Free}
\label{subsec:checkpoints}

Uma limitação crítica do Google Colab Free é o limite de sessão de aproximadamente 12 horas. Para um treino de 5 epochs com $\sim$20.500 registros (dataset híbrido), o tempo total estimado é de 10-15 horas, exigindo múltiplas sessões.

\begin{warningbox}[Estratégia de Múltiplas Sessões]
O treino completo requer 3-4 sessões do Colab. Checkpoints são salvos a cada 500 steps (~30-40 minutos) no Google Drive, permitindo retomar o treino de onde parou.
\end{warningbox}

O fluxo de múltiplas sessões é ilustrado abaixo:

\begin{figure}[H]
\centering
\begin{tikzpicture}[
    node distance=0.8cm,
    session/.style={rectangle, draw=fiap-red, fill=fiap-red!10, rounded corners, minimum width=6cm, minimum height=1.2cm, font=\small},
    checkpoint/.style={rectangle, draw=fiap-gray, fill=fiap-light, rounded corners, minimum width=4cm, minimum height=0.8cm, font=\scriptsize},
    arrow/.style={->, thick, color=fiap-dark}
]
    % Sessões
    \node[session] (s1) {Sessão 1: Steps 0 $\rightarrow$ 3500 (Epochs 1-2)};
    \node[checkpoint, below=of s1] (c1) {checkpoint-3500 salvo no Drive};

    \node[session, below=of c1] (s2) {Sessão 2: Steps 3500 $\rightarrow$ 7000 (Epochs 2-4)};
    \node[checkpoint, below=of s2] (c2) {checkpoint-7000 salvo no Drive};

    \node[session, below=of c2] (s3) {Sessão 3: Steps 7000 $\rightarrow$ FIM (Epochs 4-5)};
    \node[checkpoint, below=of s3] (c3) {Modelo final salvo};

    % Setas
    \draw[arrow] (s1) -- (c1);
    \draw[arrow] (c1) -- node[right, font=\scriptsize] {resume} (s2);
    \draw[arrow] (s2) -- (c2);
    \draw[arrow] (c2) -- node[right, font=\scriptsize] {resume} (s3);
    \draw[arrow] (s3) -- (c3);
\end{tikzpicture}
\caption{Fluxo de treinamento com múltiplas sessões no Colab Free}
\label{fig:fluxo-sessoes}
\end{figure}

O código para retomar o treino de um checkpoint:

\begin{lstlisting}[caption={Retomando treino de checkpoint existente}, label={lst:resume-checkpoint}]
import os

def find_latest_checkpoint(checkpoint_dir):
    """Encontra o checkpoint mais recente."""
    checkpoints = [d for d in os.listdir(checkpoint_dir)
                   if d.startswith("checkpoint-")]
    if not checkpoints:
        return None
    latest = sorted(checkpoints, key=lambda x: int(x.split("-")[1]))[-1]
    return os.path.join(checkpoint_dir, latest)

# Verificar e retomar
latest_checkpoint = find_latest_checkpoint(CHECKPOINT_DIR)

if latest_checkpoint:
    print(f"Retomando de: {latest_checkpoint}")
    trainer.train(resume_from_checkpoint=latest_checkpoint)
else:
    print("Iniciando treino do zero...")
    trainer.train()
\end{lstlisting}

\subsection{Execução Local}
\label{subsec:execucao-local}

Para usuários com acesso a GPUs locais (RTX 3090, 4090, A100, etc.), o notebook \texttt{03\_fine\_tuning\_local.ipynb} oferece vantagens:

\begin{itemize}
    \item \textbf{Sem limite de sessão}: Treino contínuo sem interrupções
    \item \textbf{Auto-detecção de GPU}: Ajusta batch size conforme VRAM disponível
    \item \textbf{Paths locais}: Usa diretórios do projeto diretamente
    \item \textbf{Checkpoints menos frequentes}: 1000 steps (vs 500 no Colab)
\end{itemize}

\begin{table}[H]
\centering
\caption{Configuração automática por GPU}
\label{tab:config-gpu-local}
\begin{tabular}{lccc}
\toprule
\textbf{GPU (VRAM)} & \textbf{Batch Size} & \textbf{Grad Accum} & \textbf{Batch Efetivo} \\
\midrule
24GB+ (A100, RTX 4090) & 4 & 4 & 16 \\
16-24GB (RTX 3090, A10) & 2 & 8 & 16 \\
$<$16GB & 1 & 16 & 16 \\
\bottomrule
\end{tabular}
\end{table}

\subsection{Resultados do Treinamento}
\label{subsec:resultados-treino}

O treinamento foi executado em múltiplas sessões, utilizando tanto Google Colab (GPU T4) quanto servidor local (GPU RTX A5000). A Tabela~\ref{tab:resultados-treino} apresenta os parâmetros de uma sessão completa de 5 épocas.

\begin{table}[H]
\centering
\caption{Parâmetros de treinamento (sessão de 5 épocas)}
\label{tab:resultados-treino}
\begin{tabular}{ll}
\toprule
\textbf{Parâmetro} & \textbf{Valor} \\
\midrule
Steps totais & 5.780 \\
Parâmetros treináveis & 167M (4,27\% do total) \\
Tamanho dos adapters & $\sim$640 MB \\
\bottomrule
\end{tabular}
\end{table}

\begin{warningbox}[Múltiplas Sessões de Experimentação]
O processo de fine-tuning envolveu quatro sessões de treinamento com configurações distintas de \textit{learning rate schedule}. Durante a experimentação, identificamos um fenômeno de \textbf{colapso catastrófico} após a época 1, que degradou significativamente a qualidade do modelo. A curva de loss e a análise detalhada das sessões são apresentadas no Capítulo~\ref{cap:avaliacao}.
\end{warningbox}

\begin{table}[H]
\centering
\caption{Tempo estimado de treinamento por configuração}
\label{tab:tempo-treino}
\begin{tabular}{ccccc}
\toprule
\textbf{Ambiente} & \textbf{GPU} & \textbf{Epochs} & \textbf{Tempo} & \textbf{Sessões} \\
\midrule
Colab Free & T4 (15GB) & 5 & 10-15h & 2-3 \\
Local & RTX A5000 (24GB) & 5 & $\sim$55h & 1 \\
Local & RTX 4090 (24GB) & 5 & 3-4h & 1 \\
Local & A100 (40GB) & 5 & 2-3h & 1 \\
\bottomrule
\end{tabular}
\end{table}

% ============================================================================
\section{Persistência do Modelo}
\label{sec:persistencia}

\subsection{Salvando os Adapters LoRA}

Após o treinamento, salvamos apenas os \textbf{adapters LoRA} (não o modelo completo). Isso resulta em arquivos muito menores:

\begin{table}[H]
\centering
\caption{Comparação de tamanho de arquivos salvos}
\label{tab:tamanho-modelo}
\begin{tabular}{lc}
\toprule
\textbf{Tipo de Salvamento} & \textbf{Tamanho Aproximado} \\
\midrule
Modelo completo (merged) & 15-28 GB \\
\textbf{Adapters LoRA apenas} & \textbf{~500 MB} \\
\bottomrule
\end{tabular}
\end{table}

\begin{lstlisting}[caption={Salvando o modelo fine-tuned}, label={lst:save-model}]
# Salvar adapters LoRA
SAVE_PATH = "/content/drive/MyDrive/medical-assistant-final"
model.save_pretrained(SAVE_PATH)
tokenizer.save_pretrained(SAVE_PATH)
\end{lstlisting}

\subsection{Publicação no Hugging Face Hub}
\label{subsec:hf-hub}

Os adapters LoRA ($\sim$640~MB) excedem o limite de 100~MB por arquivo do GitHub. Para garantir acesso público ao modelo treinado, os adapters são publicados no \textbf{Hugging Face Hub}, o repositório padrão da comunidade para modelos de linguagem.

Ambos os notebooks incluem uma célula dedicada ao upload:

\begin{lstlisting}[caption={Upload do modelo para o Hugging Face Hub}, label={lst:hf-upload}]
from huggingface_hub import HfApi, login

login(token=os.getenv("HF_TOKEN"))
api = HfApi()
api.create_repo(repo_id=HF_HUB_REPO, exist_ok=True)
api.upload_folder(
    folder_path=FINAL_MODEL_DIR,
    repo_id=HF_HUB_REPO,
    commit_message="Upload adapters LoRA fine-tuned",
)
\end{lstlisting}

O modelo publicado inclui um \textit{model card} (README.md) com descrição completa: hiperparâmetros, dados de treino, resultados, instruções de uso e limitações.

\subsection{Carregando o Modelo para Uso}

O modelo pode ser carregado de duas fontes --- path local ou diretamente do Hugging Face Hub:

\begin{lstlisting}[caption={Carregando o modelo fine-tuned para inferência}, label={lst:load-model}]
from transformers import AutoModelForCausalLM, AutoTokenizer

# Opcao 1: path local
model_id = "models/medical-assistant-final"
# Opcao 2: Hugging Face Hub (download automatico)
model_id = "zagari/medical-assistant-mistral-7b-ft"

tokenizer = AutoTokenizer.from_pretrained(model_id)
model = AutoModelForCausalLM.from_pretrained(
    model_id,
    torch_dtype="float16",
    device_map="auto",
)
\end{lstlisting}

Na aplicação, a variável de ambiente \texttt{FINE\_TUNED\_MODEL} controla a origem. O \texttt{.env.example} do projeto já vem configurado com o repositório do Hub, permitindo que quem clonar o projeto utilize o modelo sem precisar re-treinar.

\subsection{Exportação para GGUF (Inferência sem GPU)}
\label{subsec:gguf}

Os formatos LoRA + \texttt{bitsandbytes} 4-bit utilizados no treinamento requerem uma GPU NVIDIA com CUDA. Para permitir a execução do assistente em \textbf{máquinas sem GPU} --- como notebooks com Apple Silicon (M1/M2/M3/M4) ou desktops com apenas CPU --- o modelo pode ser exportado para o formato \textbf{GGUF}, utilizado pelo \texttt{llama.cpp}.

O processo de exportação é feito pelo script \texttt{08\_exportar\_gguf.py} no servidor com GPU:

\begin{enumerate}
    \item \textbf{Merge}: Os adapters LoRA são mesclados com o modelo base Mistral 7B, gerando um modelo completo;
    \item \textbf{Quantização}: O modelo merged é convertido para GGUF com quantização \texttt{Q4\_K\_M} (4-bit, qualidade média-alta);
    \item \textbf{Upload}: O arquivo GGUF ($\sim$4~GB) é publicado em repositório separado no Hugging Face Hub.
\end{enumerate}

\begin{table}[H]
\centering
\caption{Comparação de formatos de distribuição do modelo}
\label{tab:formatos-modelo}
\begin{tabular}{lcll}
\toprule
\textbf{Formato} & \textbf{Tamanho} & \textbf{Requisito} & \textbf{Velocidade} \\
\midrule
LoRA + base 4-bit & $\sim$640~MB + 4~GB & GPU NVIDIA (CUDA) & $\sim$2--5s/resposta \\
\textbf{GGUF Q4\_K\_M} & \textbf{$\sim$4 GB} & \textbf{CPU ou Apple Silicon} & \textbf{$\sim$5--15s/resposta} \\
Modelo merged float16 & $\sim$14~GB & 14~GB+ RAM & $\sim$2--5min/resposta \\
\bottomrule
\end{tabular}
\end{table}

A aplicação detecta automaticamente o hardware disponível: se há CUDA, carrega os adapters LoRA via \texttt{transformers}; caso contrário, carrega o GGUF via \texttt{llama-cpp-python}. Essa detecção é transparente para o usuário (detalhes na Seção~\ref{sec:interface}).

% ============================================================================
\section{Avaliação do Fine-Tuning}
\label{sec:avaliacao-ft}

Para demonstrar o \textbf{delta} (melhoria) obtido com o fine-tuning, implementamos uma avaliação comparativa rigorosa.

\subsection{Metodologia}

A avaliação segue o protocolo:

\begin{enumerate}
    \item \textbf{Baseline}: Executar 50 perguntas médicas no Mistral 7B \textit{sem} fine-tuning
    \item \textbf{Fine-tuned}: Executar as mesmas 50 perguntas no modelo \textit{após} fine-tuning
    \item \textbf{Comparação}: Análise quantitativa (métricas) e qualitativa (lado a lado)
\end{enumerate}

\subsection{Conjunto de Perguntas de Avaliação}

Criamos um conjunto de 50 perguntas médicas em português, distribuídas em 5 categorias:

\begin{table}[H]
\centering
\caption{Distribuição das perguntas de avaliação por categoria}
\label{tab:perguntas-avaliacao}
\begin{tabular}{llc}
\toprule
\textbf{Categoria} & \textbf{Exemplo de Pergunta} & \textbf{Qtd} \\
\midrule
Sintomas & ``Quais são os sintomas do diabetes tipo 2?'' & 10 \\
Tratamento & ``Como tratar uma crise aguda de asma?'' & 10 \\
Medicamentos & ``Quais são os efeitos colaterais da metformina?'' & 10 \\
Prevenção & ``Como prevenir doenças cardiovasculares?'' & 10 \\
Emergências & ``O que fazer em caso de suspeita de AVC?'' & 10 \\
\midrule
\textbf{Total} & & \textbf{50} \\
\bottomrule
\end{tabular}
\end{table}

\subsection{Métricas de Avaliação}

Utilizamos métricas automáticas complementares para avaliar a qualidade das respostas:

\begin{table}[H]
\centering
\caption{Métricas utilizadas na avaliação}
\label{tab:metricas}
\begin{tabular}{lp{8cm}}
\toprule
\textbf{Métrica} & \textbf{Descrição} \\
\midrule
\textbf{BLEU} & Mede overlap de n-gramas entre resposta gerada e referência. Valores de 0 a 1, onde maior é melhor. \\
\textbf{ROUGE-1} & Overlap de unigramas (palavras individuais). \\
\textbf{ROUGE-2} & Overlap de bigramas (pares de palavras). \\
\textbf{ROUGE-L} & Maior subsequência comum (captura estrutura). \\
\textbf{Comprimento} & Tamanho médio das respostas (em caracteres). \\
\bottomrule
\end{tabular}
\end{table}

\subsection{Scripts de Avaliação}

O pipeline de avaliação é implementado em dois scripts:

\begin{lstlisting}[style=bash, caption={Execução da avaliação}]
# 1. Avaliar modelo BASE (Mistral 7B sem fine-tuning)
python scripts/05_avaliar_baseline.py

# 2. Avaliar modelo FINE-TUNED e gerar comparacao
python scripts/06_avaliar_finetuned.py --model-path /path/to/model
\end{lstlisting}

Os resultados são salvos em:

\begin{itemize}
    \item \texttt{data/evaluation/baseline\_responses.json} --- Respostas do modelo base
    \item \texttt{data/evaluation/finetuned\_responses.json} --- Respostas do modelo fine-tuned
    \item \texttt{data/evaluation/comparacao\_resultados.json} --- Métricas comparativas
    \item \texttt{data/evaluation/comparacao\_resultados.md} --- Relatório lado a lado
\end{itemize}

\subsection{Processo Experimental}

O processo de fine-tuning não foi linear. Foram realizadas \textbf{múltiplas sessões de treinamento} com diferentes configurações de hiperparâmetros --- variando learning rate, número de epochs, e estratégias de warmup --- para encontrar a configuração ótima que evitasse problemas como \textit{catastrophic collapse} (colapso catastrófico), um fenômeno observado em alguns experimentos onde o modelo degrada significativamente após muitas épocas.

A análise completa dessas sessões experimentais, incluindo:
\begin{itemize}
    \item Comparação de 7 checkpoints de diferentes sessões de treinamento
    \item Avaliação multi-checkpoint com métricas BLEU, ROUGE e similaridade semântica
    \item Identificação do checkpoint ótimo (checkpoint-1500)
    \item Análise do fenômeno de catastrophic collapse
\end{itemize}

\noindent é apresentada em detalhes no Capítulo~\ref{cap:avaliacao}.

\begin{infobox}[Limitações da Avaliação Automática]
Métricas como BLEU e ROUGE capturam similaridade textual, mas não avaliam completamente a \textit{correção clínica} das respostas. Uma avaliação humana por profissionais de saúde seria necessária para validação completa em um cenário de produção.
\end{infobox}

% ============================================================================
\section{Resumo do Capítulo}

Este capítulo apresentou o processo completo de fine-tuning do modelo Mistral 7B para o domínio médico:

\begin{enumerate}
    \item \textbf{Modelo Base}: Mistral 7B Instruct v0.3, escolhido por ser generalista (permite demonstrar delta)
    \item \textbf{Dataset Híbrido}: $\sim$20.500 pares QA compostos por:
    \begin{itemize}
        \item 78\% MedQuAD traduzido para português ($\sim$16.400 pares)
        \item 20\% MedQuAD original em inglês ($\sim$4.100 pares)
        \item 2\% QA de medicina brasileira (55 pares sobre dengue, Chagas, SUS, etc.)
    \end{itemize}
    \item \textbf{Técnica}: QLoRA 4-bit ($r=64$, $\alpha=128$), permitindo treino em GPU T4 (15GB)
    \item \textbf{Treinamento}: 5 epochs, 5.780 steps, $\sim$55h em GPU RTX A5000 (24GB)
    \item \textbf{Publicação}: Adapters LoRA ($\sim$640~MB) publicados no Hugging Face Hub
    \item \textbf{GGUF}: Versão quantizada ($\sim$4~GB) para execução sem GPU (CPU/Apple Silicon)
    \item \textbf{Avaliação}: 50 perguntas médicas, comparação baseline vs fine-tuned
\end{enumerate}

O modelo resultante é publicado no Hugging Face Hub em dois formatos --- LoRA (para GPU) e GGUF (para CPU) --- e utilizado como o ``cérebro'' do agente de raciocínio no grafo multi-agente descrito no Capítulo~\ref{cap:multiagente}.
